\documentclass[10pt,reqno]{amsart}
\usepackage[T1]{fontenc}\usepackage{lmodern}
\usepackage[a4paper,margin=24mm]{geometry}
\usepackage{amsmath,amssymb,mathtools,microtype}
\usepackage[hidelinks]{hyperref}
\newcommand{\Z}{\mathbb Z}\newcommand{\Q}{\mathbb Q}
\newtheorem{theorem}{Theorem}\newtheorem{lemma}[theorem]{Lemma}
\title[Trace descent and cofactor capacity]{Rational trace descent and a sharp cofactor-capacity theorem}
\author{The Clankers}\date{21 September 2026}
\begin{document}
\begin{abstract}
A positive rational Erd\H{o}s--Straus candidate with integral tail sum and first
integer denominator $a=cq$, $c\in\{1,2,3,6\}$, $q$ prime and coprime to $c$,
returns an integer middle-channel solution at the same $p\equiv1\pmod{24}$.
The full target capacity $U\mid j^2$, $AU\equiv B\pmod{4j-1}$ has three persistent
branches and a finite inverse system with sharp bound
$4j-1\le\max\{B(A-1)^2-A,A(B-1)^2-B\}$. Actual prime sources require the finite
branch or have an empty complete box. Universal ES occupancy is not asserted.
\end{abstract}\maketitle

\section{Original trace candidates and a same-prime return}
Let $p\equiv1\pmod4$ be prime, $p/4<a<p/2$, $R=4a-p$, and $u\mid a^2$.
The labelled rational returns are
\[
 E:\left(a,\frac{pa+a^2/u}{R},\frac{pa+p^2u}{R}\right),\qquad
 M:\left(a,\frac{p(a+a^2/u)}R,\frac{p(a+u)}R\right).
\]
They satisfy $4/p$ exactly. Their tail sums are integral precisely when
$R\mid(4u+1)^2$ and $R\mid(a+u)^2$, respectively. The unsquared gates are the
original integer conditions. Their tail denominators are
$R/\gcd(R,4u+1)$ and $R/\gcd(R,a+u)$.
Every rational integral-tail-trace pair at the integer $a$ has one of these
markings: its positive integer factors $b=Ry-pa$, $c_1=Rz-pa$ have
$bc_1=(pa)^2$, and their $p$-valuations distinguish E from M. This is the
factor-pair framework of \cite[\S2]{ET}, with the relaxed trace gate retained.

\begin{theorem}\label{t:core}
Suppose $p\equiv1\pmod{24}$, $a=cq$, $q$ prime, $\gcd(c,q)=1$, and
$c\in\{1,2,3,6\}$. Every rational trace candidate above produces an original
integer M state at $p$. If its input tails are nonintegral, the returned middle
cofactor is strictly smaller than $R$.
\end{theorem}
\begin{proof}
For every prime $\ell\mid a$, reciprocity and $R\equiv-p\pmod\ell$ give
$(\ell/R)=(\ell/p)$; the supplementary law handles two. The squared gates
therefore imply $(u/p)=-1$ in E and $(u/p)=-(a/p)$ in M. All factors of $c$
are residues, so $q$ is a nonresidue. Hence, for some actual $d\mid c^2$,
\[
 E:u=qd;\qquad M:u=d\text{ or }a^2/d.
\]
Write $R=Qv^2$, where $Q$ is its squarefree part, and set $j=(Q+1)/4$.
Then $Q\equiv3\pmod4$, $3\le Q<p$. Since $4c\mid24$, every unit square
modulo $4c$ is one. The original equation $Qv^2=4cq-p$ gives $c\mid j$.
Modulo $Q$, the E gate is $dp\equiv-c$, and the M gate is $p\equiv-4d$.
Use $U=cj/d$ in E and $U=d$ in M. Writing $j=cm$ proves $U\mid j^2$;
$4j\equiv1\pmod Q$ proves $Q\mid p+4U$.

For completeness the exact integer return is
\begin{align*}
 R'&=(p+4U)/Q,&a'&=(p+R')/4=jR'-U,\\
 g'&=\gcd(j,U),&(h',r',\lambda')&=(g'^2/U,U/g',j/g'),\\
 s'&=R'\lambda'-r',&u'&=U.
\end{align*}
Prime valuations give positive integral $h',r',\lambda'$, with
$j=h'r'\lambda'$, $U=h'r'^2$ and $\gcd(r',\lambda')=1$.
Since $4U\le(Q+1)^2/4<p(Q-1)$, $0<R'<p$ and $p/4<a'<p/2$.
The identity $a'=h'r's'$ makes $s'>0$. A prime common to $r',R'$ would divide
$p$, impossible since $r'\le j<p$. Thus the source is normalized and
$r'+s'=R'\lambda'$. Its ordered denominators are
\[
 (a',ph's'\lambda',ph'r'\lambda'),
\]
with inverse $U=pa'^2/(R'y-pa')$. Nonintegral tails require a nonsquarefree
$R$, so $Q<R$. No decrease of $R'$ is asserted.
\end{proof}
For a general residue core, the same proof works on the exact domain $c\mid j$.
Its equivalent original condition is $v^2\equiv p\pmod{4c}$, not just a statement
about the subgroup generated by the factors.

The full inverse fibre retains $c$, the E/M tag and M orientation. Read $Q,j,U$
from the target; set $d=U$ in M or $d=cj/U$ in E. Choose odd $v$ with $Qv^2<p$,
require $q=(p+Qv^2)/(4c)$ prime and coprime to $c$, then check $d\mid c^2$ and
\emph{the full original squared gate}. These finite tests recover every input.
At fixed $p,Q,c$, an E and an M word have the same target precisely when
$d_Ed_M=cj$. Their integer multiplicities remain separate until pushforward.

\section{Complete capacity at a rational target}
The two channels require respectively $U\equiv d$ and $4dU\equiv c$ at $Q$.
This makes the following two-parameter problem unavoidable. The $A=1$ case is
an inherited workbench theorem \cite[eqs. (7)--(12)]{Capacity}.

\begin{theorem}\label{t:capacity}
Let $A,B,j$ be positive integers and $Q=4j-1>\max(A,B)$. All words
$U\mid j^2$ satisfying $AU\equiv B\pmod Q$ are the disjoint union of
\[
 U=B/A,\qquad U=4Bj/A,\qquad U=16Bj^2/A,
\]
when each is an integer divisor of $j^2$, and the finite records
\begin{gather*}
 n,k>0,\quad nk<AB,\quad
 Q=\frac{Bk+An-2AB}{AB-nk}>\max(A,B),\quad Q\equiv3\pmod4,\\
 j=(Q+1)/4,\quad U=(B+nQ)/A\in\Z_{>0},\quad
 V=(A+kQ)/(16B)\in\Z_{>0}.
\end{gather*}
Every finite record satisfies
\[
 \boxed{Q\le H(A,B):=\max\{B(A-1)^2-A,A(B-1)^2-B\}.}
\]
The inverse is $V=j^2/U$, $n=(AU-B)/Q$, $k=(16BV-A)/Q$.
\end{theorem}
\begin{proof}
Multiplying $AU\equiv B$ by $16V$, with $UV=j^2$, proves $16BV\equiv A$.
The strict $Q$-bound makes $n,k$ nonnegative, and expansion gives
\[
 (AB-nk)Q=Bk+An-2AB.\tag{*}
\]
The cases $n=0$ and $k=0$ are the constant and quadratic branches. If $n,k>0$
and $nk>AB$, the left side is negative while
$Bk+An\ge2\sqrt{ABnk}>2AB$. If $nk=AB$, equality forces $(n,k)=(B,A)$,
the linear branch. Every remaining word has the stated finite record, and
conversely (*) gives $UV=j^2$. The inverse proves disjointness.

For a finite record let $z=nk\le AB-1$. Convexity on $1\le n\le z$ gives
$An+Bz/n\le\max(Az+B,A+Bz)$. Hence $Q$ is bounded by the maximum of
\[
 \frac{Az+B-2AB}{AB-z},\qquad\frac{A+Bz-2AB}{AB-z}.
\]
These increase in $z$, with derivative numerators $B(A-1)^2$ and $A(B-1)^2$.
Putting $z=AB-1$ proves the result; $AB=1$ has no finite branch.
\end{proof}
There are at most $(AB-1)(1+\log(AB-1))$ finite quotient pairs when $AB>1$.
The omitted domain $Q\le\max(A,B)$ is a literal finite $j$-range. No rational
word is rounded, and no generating subgroup replaces $U\mid j^2$.

The bound is sharp. For $\ell\ge1$, $A=4\ell+1$ and $2\le B<A$, set
\[
 Q=B(A-1)^2-A,\quad j=\ell(4B\ell-1),\quad
 U=(4B\ell-1)^2,\quad V=\ell^2.
\]
Then $n=AB-1,k=1$, and $Q=H(A,B)$. This integer family does not assert
infinitely many prime ES parameters.

\section{Original examples and the remaining boundary}
At $p=1009$, the rational M candidate
\[
 (a,y,z)=(321,215926/5,6054/5),\qquad a=3\cdot107,
\]
has an empty original shell. Theorem~\ref{t:core} returns
\[
 4/1009=1/276+1/92828+1/3027.
\]
This goes beyond prime distinguished denominators at the same $p$.

The finite branch is needed at the certified hard prime $p=2377702849$:
\[
 a=288\cdot2064031=594440928,\quad R=60863=503\cdot11^2,
 \quad u=81\cdot2064031.
\]
The E tail quotient is $3867581/11$, and its tail sum is an integer.
Its reduced capacity is $9U\equiv8\pmod{503}$ at $j=126$. The complete box
has exactly one target word, $U=3969$, with $(V,n,k)=(4,71,1)$ and
$503=H(9,8)$. The constant word is nonintegral; the linear word $448$ is not a
divisor of $126^2$; the quadratic branch is unavailable. The ordered return is
\[
 (595607481,44958019189187726,299590558974).
\]

For both channels the extension to $c=4$ fails at a specific source:
\[
 p=9601,\ a=4\cdot739=2956,\ R=2223,\ u=8,\qquad
 (y,z)=(14190278/3,38404/3).
\]
The old shell is empty in E and M. Its squarefree cofactor is $Q=247$, $j=62$.
The complete divisors of $j^2$ are
$1,2,4,31,62,124,961,1922,3844$; none is $8\pmod{247}$. Theorem~\ref{t:capacity}
also certifies the failure: $247>H(1,8)=41$, and none of the three persistent
words is available. Nevertheless
\[
 4/9601=1/2405+1/46180810+1/1248130.
\]
Thus the failure concerns an exact source-changing map, not ES.

The full workbench proof also handles all E candidates at $a=2^e q$, $e\le4$,
by three available binary-exponent intervals, and supplies every inverse fibre.
The standard-library replay tests all first-half divisors for primes through
3000 and all capacities $A,B\le24$, $j\le1500$, with a separate implementation.
Finite checks do not prove universal occupancy. The next original existence
obligation has not been assumed, and no historical-priority claim is made.

\begin{thebibliography}{3}
\bibitem{ET} C. Elsholtz and T. Tao (2013).
Counting the number of solutions to the Erd\H{o}s--Straus equation on unit fractions.
\emph{J. Aust. Math. Soc.} 94(1), 50--105. arXiv:1107.1010, \S2.
\bibitem{Capacity} The Clankers (2026, September 20).
\emph{Exact cofactor capacity and fixed-target returns}.
Erd\H{o}s--Straus Foundation, \texttt{es-turn07-capacity-completion-20260920},
\texttt{core.tex}, equations (7)--(12), at revision
\texttt{49976b17491ba927df8bc7f2c2f81d39dfb15ce8}.
\bibitem{Prime} The Clankers (2026, September 21).
\emph{Prime-shell trace descent and an exact global denominator cutoff}.
Preserved source \texttt{proof.md}, complete prime-shell classification and return.
\end{thebibliography}
\end{document}
